% Propuesta separada; no modifica el bloque común del artículo I.
% Procedencia: artículo II, revisión 05, registro_incidencias.tex y
% registro_imagen_integral.tex; aclaraciones de índices de revisiones 06/07.
% Inserción prevista inmediatamente antes de subsec:k-terminal.
\subsection{Évaluation des incidences et détermination des coordonnées transversales}
\label{iii:r25:incidencias}

Le registre dodécaphasique conserve une histoire de visites et de traces
orientées. Ses coordonnées transversales s’obtiennent au moyen d’une
composition d’opérations : évaluation arithmétique des visites,
dénombrement par fenêtres, normalisation décimale et différences cycliques.
La reconstruction du registre à partir de ces différences constitue
une autre opération, dont l’existence et l’unicité seront démontrées ensuite.
Cette séparation permet de préciser tant les données utilisées par
l’extracteur que l’information qu’il transmet à la lecture harmonique.

Soit $\mathscr L$ un registre d’incidences produit par la sélection,
le transport et la mise à jour d’une histoire APP--TRIT--TPK.
Chaque visite conserve son chemin, sa position APP, son feuillet, son instant,
sa multiplicité et sa classe de frontière ; chaque trace conserve ses membres,
son orientation, sa multiplicité et sa longueur dans l’unité déclarée par son
lecteur. Les fenêtres sont
\[
 E_m=\{9(m-1)+1,\ldots,9m\},\qquad 1\leq m\leq12.
\]
L’évaluation du feuillet de somme ou de produit détermine l’entier
$n(v)$ de chaque visite et sa division positive
\[
 n(v)=r(v)+9q(v),\qquad 1\leq r(v)\leq9.
\]
Le résidu se calcule sur cette évaluation. La classe de frontière d’une
visite de résidu $9$ est conservée au moyen de
$\epsilon_\partial(v)=1$ dans la couronne et de $-1$ à l’intérieur.

La présentation d’incidence du registre utilise les trois dénombrements
\begin{align}
 A_m&=\sum_{v\in E_m}w(v)
       \mathbf1_{\{1,3,5,7\}}(r(v)),\label{iii:r25:recuento-a}\\
 C_m&=\sum_{j\geq0}
       \bigl(\nu^-_{120,m}(j)-\nu^+_{120,m}(j)\bigr),
       \label{iii:r25:recuento-c}\\
 V_m&=\sum_{v\in E_m}w(v)
       \mathbf1_{\{9\}}(r(v))\epsilon_\partial(v).
       \label{iii:r25:recuento-v}
\end{align}
Ici, $w(v)$ est la multiplicité d’incidence et
$\nu^\pm_{120,m}(j)$ compte les traces orientées du canal indiqué
de longueur réduite $120(1+3j)$, selon l’énumération du registre.
La famille de traces et leur multiplicité sont des données de l’évaluation,
avec la règle de production qui leur correspond. La somme entière nécessite
un support fini ou une construction qui détermine son sens.
La longueur réduite d’une trace et le nombre d’instants de la
fenêtre sont des grandeurs typées différentes ; l’unité du lecteur
spécifie leur relation. Les symboles $A_m,C_m,V_m$ désignent ici
des dénombrements entiers, distincts de la paire angulaire et de l’opérateur constitutif.

Soit $[x]_{10}\in\{0,\ldots,9\}$ le représentant euclidien.
La normalisation conserve également le quotient dans
\[
 x=[x]_{10}+10\lfloor x/10\rfloor.
\]
On forme le bloc
\begin{equation}
 z_m=100[A_m]_{10}+10[C_m]_{10}+[V_m]_{10},
 \qquad 0\leq z_m\leq999.
 \label{iii:r25:bloque}
\end{equation}
Avec des indices cycliques modulo $12$, soient
\[
 (D_a z)_m=z_m-z_{m+a},\qquad a\in\{3,4\},
 \qquad Q(z)=\sum_{m=1}^{12}z_m.
\]
La composition d’incidence s’exprime par
\begin{equation}
 \mathcal E_{\rm cnt}(\mathscr L)
 =\mathcal D(z):=(D_3z,D_4z,Q(z))\in\mathbb Z^{25}.
 \label{iii:r25:composicion}
\end{equation}
Ses arguments sont les visites et les traces avec leurs règles de dénombrement.
L’évaluation ne reçoit ni $K$, ni $U$, ni $\alpha$, ni les coordonnées
transversales attendues. Pour l’identifier à
$\mathcal E_{108}^{90,120}$, on conserve en outre la composition avec
la famille et les incidences effectives de l’état terminal.
L’inversibilité de $\mathcal D$ détermine une reconstruction ultérieure ;
à elle seule, elle ne sélectionne pas cette famille d’incidences.

\begin{proposition}[Information résiduelle déterminée par l’extracteur]
\label{iii:r25:residuos}
Soient $N,N'\in\mathbb Z^{12\times3}$ deux registres de dénombrements,
ordonnés par $(A,C,V)$. Pour les opérations
\eqref{iii:r25:bloque}--\eqref{iii:r25:composicion},
\[
 \mathcal E_{\rm cnt}(N)=\mathcal E_{\rm cnt}(N')
 \quad\Longleftrightarrow\quad N-N'\in10\mathbb Z^{36}.
\]
Les $36$ résidus sectoriels déterminent exactement la sortie de
$25$ coordonnées. Les quotients et les provenances constituent
une information supplémentaire du registre.
\end{proposition}
\begin{proof}
L’égalité modulo $10$ produit les mêmes blocs et leurs mêmes
coordonnées transversales. Réciproquement, l’égalité des
coordonnées donne $D_3(z-z')=D_4(z-z')=0$. Les pas $3$ et $4$
engendrent le cycle de longueur $12$, donc $z-z'$ est constant.
L’égalité de $Q$ annule cette constante. Le développement de chaque entier
de $0$ à $999$ en trois chiffres décimaux est unique ; on retrouve ainsi
les trois résidus de chaque fenêtre.
\end{proof}

La proposition est valable sur tout le domaine entier indiqué. Elle ne
postule pas que les $36$ classes soient indépendantes sur le sous-ensemble
des registres admissibles produit par TPK.

\subsection{Image entière et inversion des coordonnées transversales}
\label{iii:r25:imagen-seccion}

Pour exprimer l’inversion, on réindexe les blocs par
$i\in\mathbb Z/12\mathbb Z$, de $0$ à $11$. Soient
$(Sz)_i=z_{i+1}$, $D_3=I-S^3$ et $D_4=I-S^4$.
Étant donné un triplet entier $(b,c,q)$, définissons
\begin{equation}
 w_i=c_i-b_{i+1},\qquad
 P_0=0,\quad P_i=\sum_{j=0}^{i-1}w_j\ (1\leq i\leq12),
 \qquad T(w)=\sum_{i=0}^{11}P_i.
 \label{iii:r25:incrementos}
\end{equation}

\begin{theorem}[Caractérisation de l’image entière]
\label{iii:r25:imagen}
Un triplet $(b,c,q)\in\mathbb Z^{25}$ appartient à
$\operatorname{im}\mathcal D$ si et seulement si
\begin{align}
 b_i&=w_i+w_{i+1}+w_{i+2},\qquad 0\leq i<12,
       \label{iii:r25:relaciones}\\
 \sum_{i=0}^{11}w_i&=0,\label{iii:r25:cierre}\\
 q+T(w)&\equiv0\pmod {12}.\label{iii:r25:congruencia}
\end{align}
Son unique préimage entière est
\begin{equation}
 z_i=\frac{q+T(w)}{12}-P_i,\qquad0\leq i<12.
 \label{iii:r25:inversa}
\end{equation}
Les douze relations \eqref{iii:r25:relaciones} et la relation
\eqref{iii:r25:cierre} sont linéairement indépendantes sur
$\mathbb Q$.
\end{theorem}
\begin{proof}
Pour une sortie de $\mathcal D$,
\[
 c_i-b_{i+1}
 =(z_i-z_{i+4})-(z_{i+1}-z_{i+4})=z_i-z_{i+1}.
\]
Trois accroissements consécutifs ont pour somme $b_i$ et leur somme cyclique
totale est nulle. De plus, $P_i=z_0-z_i$, d’où
$q=12z_0-T(w)$. On obtient les conditions et la formule inverse.
Réciproquement, la congruence rend entières les coordonnées de
\eqref{iii:r25:inversa} et la clôture permet de les prolonger
cycliquement. Leurs différences adjacentes sont $w_i$, donc $D_3z=b$
et
\[
 (D_4z)_i=w_i+w_{i+1}+w_{i+2}+w_{i+3}
 =w_i+b_{i+1}=c_i.
\]
La somme donne $Q(z)=q$. La formule détermine également l’unicité.
Enfin, le changement entier inversible
$(b,c)\mapsto(b,w=c-Sb)$ transforme les douze premières relations
en $b=(I+S+S^2)w$ : chacune isole une coordonnée différente de
$b$. La clôture $\sum_iw_i=0$ ajoute une relation indépendante.
\end{proof}

L’image rationnelle a pour dimension $12$. La condition entière ajoute
la congruence d’ordre $12$ de \eqref{iii:r25:congruencia}.
Les onze accroissements $w_0,\ldots,w_{10}$ et la charge $q$ sont
des coordonnées suffisantes, avec $w_{11}=-\sum_{i=0}^{10}w_i$.
Les $25$ coordonnées conservent ces $12$ degrés de liberté
au moyen de $13$ relations linéaires. Si le triplet provient en outre de
\eqref{iii:r25:bloque}, l’inverse appartient à
$\{0,\ldots,999\}^{12}$ et retrouve ses trois résidus par fenêtre.

\subsubsection{Concordance avec la transformation harmonique}

La transformation $\mathcal H_{12}$ agit au moyen de $H_4$ sur les
orbites $(r,r+3,r+6,r+9)$, $r=0,1,2$, où
\[
 H_4=\begin{pmatrix}
 1&1&1&1\\1&1&-1&-1\\1&-1&1&-1\\1&-1&-1&1
 \end{pmatrix},\qquad H_4^2=4I_4.
\]
Pour un triplet compatible, soient
\[
 B_r=\sum_{j=0}^3c_{r+3j},\qquad
 a_0=\frac{q+2B_0+B_1}{3},\quad
 a_1=a_0-B_0,\quad a_2=a_1-B_1.
\]
Le lecteur harmonique a pour blocs
\[
 (\Pi_H(b,c,q))^{(r)}
 =(a_r,b_{r+3}-b_{r+9},b_r+b_{r+6},b_r-b_{r+6}).
\]

\begin{proposition}[Concordance des reconstructions]
\label{iii:r25:hadamard}
Sur $\operatorname{im}\mathcal D$, on a
\[
 \Pi_H\mathcal D=\mathcal H_{12},\qquad
 \mathcal D^{-1}=\tfrac14\mathcal H_{12}\Pi_H.
\]
L’image de $\Pi_H$ restreinte à ce domaine est
\[
 \mathcal L_H=
 \{u\in\mathbb Z^{12}:\mathcal H_{12}u\equiv0\pmod4\}.
\]
\end{proposition}
\begin{proof}
Pour $z=\mathcal D^{-1}(b,c,q)$, soient
$h_r=\sum_{j=0}^3z_{r+3j}$. Le décalage de quatre positions
envoie chaque orbite sur la suivante ; par conséquent, $B_r=h_r-h_{r+1}$.
L’identité $q=h_0+h_1+h_2$ donne $a_r=h_r$.
Sur une orbite, écrivons $(x_0,x_1,x_2,x_3)$ pour ses coordonnées
et $d_j=x_j-x_{j+1}$. On vérifie
\[
 \begin{aligned}
 d_1-d_3&=x_0+x_1-x_2-x_3,\\
 d_0+d_2&=x_0-x_1+x_2-x_3,\\
 d_0-d_2&=x_0-x_1-x_2+x_3.
 \end{aligned}
\]
Ce sont les trois lignes restantes de $H_4x$. L’inversion et la
caractérisation entière découlent de $\mathcal H_{12}^2=4I$.
\end{proof}

La condition de congruence harmonique contrôle la sortie de $\Pi_H$ ;
l’appartenance des canaux à $\operatorname{im}\mathcal D$
nécessite en outre \eqref{iii:r25:relaciones}--\eqref{iii:r25:congruencia}.
Un contrôle négatif exact consiste à prendre $b=0$, $q=0$ et
$c=e_0-e_3$. Les trois sommes orbitales de $c$ s’annulent, de sorte
que $\Pi_H(b,c,q)=0$, mais les relations d’image ne sont pas satisfaites.
La lecture harmonique ne remplace pas le contrôle de compatibilité des
canaux.

\subsection{Conjugaison, covariance et mémoire du registre d’incidence}
\label{iii:r25:transporte}

Soit $\rho(m)=13-m$. Si la conjugaison des incidences inverse
l’ordre des fenêtres, conserve les classes arithmétique et de frontière
des visites et échange les canaux des traces, alors
\[
 (A_m^\dagger,C_m^\dagger,V_m^\dagger)
 =(A_{\rho(m)},-C_{\rho(m)},V_{\rho(m)}).
\]
Avec $c_m=[C_m]_{10}$, la normalisation décimale donne
\begin{equation}
 z_m^\dagger=z_{\rho(m)}
 +100\mathbf1_{\{c_{\rho(m)}\ne0\}}-20c_{\rho(m)}.
 \label{iii:r25:conjugacion}
\end{equation}
En effet, $[-C]_{10}=0$ si $[C]_{10}=0$ et vaut
$10-[C]_{10}$ dans les autres cas. La négation du canal
s’applique avant de reconstruire le bloc décimal.

L’application à une involution concrète conserve également les
extrémités des fenêtres. Pour un chemin
$\gamma=(x_0,\ldots,x_n)$ lu avant l’avancée,
\begin{equation}
 \sum_{k=0}^{n-1}f(x_{n-k})
 =\sum_{k=0}^{n-1}f(x_k)+f(x_n)-f(x_0).
 \label{iii:r25:extremos}
\end{equation}
Les deux sommes contiennent les mêmes termes intérieurs ; la différence
est l’extrémité finale moins l’initiale. Sur les chemins fermés, cette
correction s’annule. Dans les fenêtres ouvertes, elle est incorporée avant de
réduire modulo $10$.

\begin{proposition}[Conséquence d’une réduction périodique]
\label{iii:r25:periodo}
Si l’observable d’une famille fixe de chemins satisfait
$o(t+54)=o(t)$, si ses multiplicités sont conservées sous ce retour
et si le même lecteur local agit dans des fenêtres de neuf instants,
alors
\[
 (A,C,V)_{m+6}=(A,C,V)_m,\qquad z_{m+6}=z_m.
\]
\end{proposition}
\begin{proof}
La translation $t\mapsto t+54$ envoie chaque fenêtre bijectivement
sur celle située six positions plus loin et conserve chaque incidence
comptée et sa multiplicité. L’égalité se transmet aux dénombrements
et à leurs normalisations décimales.
\end{proof}

Une présentation à douze blocs distincts conserve donc
une information qui ne se factorise pas par cette observable de période $54$ :
sélection des incidences, orientation, extrémités, multiplicité ou
mémoire. La proposition délimite l’effet de cette réduction ; elle
n’attribue pas une période $54$ à l’histoire complète.

\subsubsection{Détermination sur les orbites et accumulation d’événements}

Fixer uniquement $Q(z)=q$ laisse une classe affine du réseau
\[
 \mathcal Z_0=\{v\in\mathbb Z^{12}:\sum_i v_i=0\}
 =\bigoplus_{i=0}^{10}\mathbb Z(e_i-e_{11}).
\]
Par exemple, $100\mathbf1+e_0$ et $100\mathbf1+e_1$ ont pour
charge $1201$ et satisfont les contrôles d’image au moyen de leurs
canaux respectifs. Ce sont des témoins des équations de compatibilité,
sans attribution d’appartenance aux histoires terminales HMT.

\begin{proposition}[Covariance cyclique]
\label{iii:r25:covariancia}
Soit $\rho$ la translation des fenêtres d’un registre $R$, de
période minimale $p\mid12$, et soit $S$ la translation de ses
coordonnées. Un lecteur covariant sur cette orbite est déterminé
par $z\in\operatorname{Fix}(S^p)$. La condition $Q(z)=q$
admet des solutions entières si et seulement si $12/p$ divise $q$ ;
lorsqu’elles existent, elles forment une classe affine de rang $p-1$.
\end{proposition}
\begin{proof}
La règle $F(\rho^jR)=S^jz$ est bien définie exactement si
$S^pz=z$. Ces vecteurs répètent $12/p$ fois un bloc de
longueur $p$, et leur charge est $\frac{12}{p}\sum_{i=0}^{p-1}z_i$.
La formule prouve la divisibilité et le rang. La covariance des
canaux découle de $D_aS=SD_a$ et $QS=Q$.
\end{proof}

Dans la coordinatisation harmonique, l’action correspondante est
$\widehat S_H=\mathcal H_{12}S\mathcal H_{12}/4$.
Si l’origine est marquée, l’action transporte également cette marque.
La période de l’orbite indiquée ici ne s’identifie pas à celle de
l’histoire enrichie. Pour une lecture initiale $F_{108}$ déjà
définie, la naturalité sous troncature s’exprime au moyen de
$F_N=F_{108}\circ\tau_{N,108}$ ; cette identité conserve
l’évaluation initiale sans sélectionner sa règle de production.

\begin{proposition}[Composition additive des contributions d’incidence]
\label{iii:r25:eventos}
Supposons que l’état initial et chaque événement $a_t$ produisent,
au moyen de règles préalablement fixées sur le registre, des contributions
$(\delta w_t,\delta q_t)$ telles que
\[
 \sum_i\delta w_{t,i}=0,
 \qquad \delta q_t+T(\delta w_t)\equiv0\pmod {12}.
\]
Avec des conditions initiales compatibles $(w^{(0)},q^{(0)})$,
la mise à jour par somme détermine un unique registre à chaque étape.
Son accroissement est
\[
 \delta z_{t,i}=
 \frac{\delta q_t+T(\delta w_t)}{12}-P_i(\delta w_t).
\]
La construction respecte la concaténation des segments et conserve
la lecture d’un préfixe sous les prolongements ultérieurs.
\end{proposition}
\begin{proof}
Les conditions de compatibilité se conservent par somme et
l’inverse \eqref{iii:r25:inversa} est additif en $(w,q)$.
La récurrence fixe chaque état. La somme sur des segments concaténés
est la somme de leurs contributions ; un préfixe reçoit les mêmes
termes avant et après le prolongement.
\end{proof}

Cette proposition donne un critère suffisant pour l’accumulation
additive, non une exigence de linéarité de tout lecteur TPK.
Lorsque les événements se transportent par des actions non triviales, leurs
contributions s’expriment dans la coordinatisation terminale au moyen du cocycle
affine correspondant. Les règles d’événement et l’état initial
conservent dans les deux cas leur provenance d’incidence.

Le développement précédent détermine le passage
\[
 \mathscr L\longmapsto(A,C,V)\longmapsto z
 \longmapsto(D_3z,D_4z,Q(z))\longmapsto\mathcal H_{12}z
\]
avec les conditions de chaque application. Son domaine est distinct du
registre régional de l’appendice~\ref{iii:nucleo-finito} : les
équations $X_aL_a=Y_a$ y déterminent les relèvements à partir de
transitions ternaires. Ici, les incidences des visites et des traces
déterminent des blocs entiers et leurs canaux. L’unicité d’une
reconstruction ne remplace pas la production du registre qu’utilise
l’autre. Les deux conservent leurs primitives, opérations et preuves dans
la chaîne APP--TRIT--TPK.
